\chapter{系统需求分析}

\section{系统建设目标}
本系统面向人工智能专业的医学图像处理实验，目标是搭建一个可操作的 MRI 医学图像拟生成平台。平台需要把数据预处理、StyleGAN2 二维生成、连续切片组织、点云初始化、3DGS 训练结果展示和 AI 辅助说明放在同一工作流程中。相比只使用命令行脚本，平台更强调操作路径清楚、结果可以直接查看，实验过程也便于再次复核。

使用对象主要包括三类：进行医学图像生成实验的学生，学习 StyleGAN2 与 3DGS 的学习者，以及需要展示生成式医学影像流程的教学人员。平台不面向临床诊断，也不输出医学诊断结论，其定位是算法教学、毕业设计展示和工程实验验证。

\section{功能需求分析}
功能设计围绕“从 MRI 数据到二维生成，再到三维展示”这一主线展开。平台需要提供数据预处理、模型训练、切片生成、连续序列生成、点云初始化、3DGS 结果管理和 AI 辅助问答等功能。各模块不是彼此独立的展示项，而是按照实验流程前后衔接。

数据预处理是整条流程的起点。平台读取 NIfTI 格式 MRI 文件后，自动提取轴向切片，完成灰度归一化、尺寸缩放、质量过滤和 PNG 保存，再调用 StyleGAN2 数据集工具生成训练压缩包。为了让用户判断处理是否正常，平台需要返回处理日志，说明输入路径、输出路径、目标尺寸、有效切片数量和异常文件信息。

模型训练功能承担二维生成模型的构建任务。用户可以设置学习率、批大小和训练迭代量等基础参数。由于 StyleGAN2 训练耗时较长，训练任务需要在后台运行。系统保存模型快照、配置文件和训练日志，为论文复核提供实验材料。

切片生成功能把训练权重转化为可观察结果。平台加载训练得到的 \texttt{.pkl} 权重文件，根据随机种子、生成数量和截断参数输出二维 MRI 切片。生成结果保存到结果目录，同时直接在前端显示。该功能可用于论文插图整理和答辩演示。

连续序列生成用于连接二维样本和三维表达。系统在两个随机种子对应的潜变量之间插值，得到连续变化的切片序列，并把独立二维样本组织为近似体数据结构的序列。该序列为点云初始化和 3DGS 训练提供输入。

点云初始化是二维结果进入三维空间的接口。系统读取连续切片目录，将非背景像素映射为空间点，并把灰度值转换为颜色信息，最终输出 PLY 格式点云。该模块使后续 3DGS 训练具有明确的空间输入。

3DGS 结果管理更关注实验产物能否被追踪和复核。平台需要归档训练日志、相机参数、checkpoint、渲染图像和训练后点云，并通过路径、日志和图像预览展示三维训练结果。三维模块不只是生成文件，还需要留下可以解释和检查的实验材料。

AI 辅助问答服务于教学说明。它用于解释 MRI 图像生成、StyleGAN2 参数、3DGS 基本概念和平台操作流程，但不参与核心训练和推理。也就是说，即使外部模型接口不可用，平台主要功能仍应保持可运行。

\section{性能与可用性需求分析}
性能需求主要体现在两类场景。MRI 批量预处理、StyleGAN2 训练和 3DGS 训练耗时较长，需要交给后端任务进程执行，避免阻塞 React 前端界面。切片生成、图像预览、影像分析和点云初始化更接近日常演示操作，平台应尽量在可接受时间内返回日志、图像或输出路径。

可用性要求集中在低门槛操作和清晰反馈上。用户不应直接面对复杂命令行参数，而应通过输入框、滑块、按钮和日志区完成操作。遇到数据路径不存在、模型权重缺失或生成数量异常等问题时，平台需要返回明确提示；正常运行时，则展示图像画廊、输出路径和关键日志。

系统扩展性需求主要体现在模型、三维重建和评价指标三个方向。二维生成部分需要支持替换不同 StyleGAN2 权重或不同训练数据。三维部分需要从点云初始化扩展到更完整的 3DGS 训练与渲染。评价部分则可继续加入 SSIM、FID、LPIPS、专家评价或下游任务指标。平台代码应保持模块化，避免新增功能时重写整体界面。

\section{可行性分析}
从技术可行性看，本文所用技术都有较成熟的开源实现或 Python、TypeScript 生态支持。StyleGAN2-ADA 有可复用的 PyTorch 工程，NIfTI 数据可由 nibabel 读取，图像处理可使用 NumPy 和 PIL 完成。FastAPI 可以组织后端接口和 WebSocket 通信，React 可以构建复杂前端页面，Vite 可以支撑本地开发和前端构建。当前项目已经形成 StyleGAN2 权重、生成切片、初始点云、3DGS 训练日志和渲染结果，这说明核心流程已经具备实现基础。

从数据可行性看，IXI 数据集公开提供多种脑部 MRI 数据，类型包括 T1、T2、PD、MRA 和 DTI 等。这些数据以 NIfTI 格式发布，适合教学和算法实验\cite{ixiDataset}。本文使用 T1 脑部 MRI 作为主要实验数据，能够支撑二维切片生成和三维序列组织。

从工程可行性看，系统采用“算法脚本 + FastAPI 后端 + React 前端”的组织方式。底层脚本负责计算任务，后端负责接口封装、任务调度和结果归档，前端负责参数输入、状态展示、影像查看和三维交互。三者通过 HTTP 接口、WebSocket、文件路径和日志连接。这样的组织方式更接近真实 Web 平台结构，也便于继续扩展页面和接口。

从应用可行性看，平台定位为教学与研究演示系统，并不是临床诊断软件。平台不承担医疗器械级别的验证要求。论文对生成结果的描述以算法展示、工程闭环和可视化验证为主，不把平台输出解释为临床可用数据。
